\documentclass[11pt]{article}
\usepackage{amsmath,amssymb,amsthm}
\newtheorem{lemma}{Lemma}
\title{The Biased Butterfly Dual and a One-Level Estimate}
\author{}\date{}\begin{document}\maketitle

\paragraph{Status.}
This is a complete written audit of several exact identities and one
quantitative necessary condition on an infeasibility certificate. It
does not exclude all certificates, prove a density theorem, or give a
new asymptotic HeadComplexity bound. The gap after the estimate is stated
explicitly below. The companion checker verifies the finite algebra,
not the missing induction.

\paragraph{Notation.}
There are $H=2^d$ wires and butterfly layers $k=0,\ldots,d-1$.
At layer $k$, a block $g$ fixes the input bits above $k$; its queries
are indexed by lower output prefixes $s<2^k$. A switch word
$\theta_{k,g,s}$ determines the routed permutation $\pi$.
Write $\sigma_{k,g,s}=1-2\theta_{k,g,s}$ and
\[
 u_{k,g}(h)=\mathbf1_{\{\lfloor h/2^{k+1}\rfloor=g\}}
             (1-2h_k).
\]
Choose positive digit sizes $\eta_0,\ldots,\eta_{d-1}$ and put
$p_y=\sum_j y_j\eta_j$, $P=\sum_j\eta_j<1/2$,
$b=1-2P>0$. The query time for prefix $s$ is $1-p_s$.
For residues $c_y$ and an independent bias $\beta_{k,g}$ at each block,
the desired inequalities are
\[
 \sigma_{k,g,s}\left(\beta_{k,g}+
 \sum_y u_{k,g}(\pi^{-1}y)\frac{c_y}{1-p_s-p_y}\right)>0.
 \tag{1}
\]
All residues and biases here are unrestricted real numbers.

\paragraph{Exact dual.}
The strict homogeneous system (1) is infeasible if and only if there is
a nonzero family $\lambda_{k,g,s}\geq0$ such that, with
$v_{k,g,s}=\sigma_{k,g,s}\lambda_{k,g,s}$ and
\[
 R_{k,g}(x)=\sum_s\frac{v_{k,g,s}}{1-p_s-x},
\]
one has
\[
 \sum_s v_{k,g,s}=0,\qquad
 \sum_{k,g}u_{k,g}(h)R_{k,g}(p_{\pi(h)})=0
       \quad\hbox{for every input }h.                 \tag{2}
\]
This is the strict linear alternative applied to the variables
$(c,\beta)$. The first equations are exactly the bias coordinates.
In particular every level-zero dual block is zero, since it has one
query. A constant-sign target block also has zero dual mass.

\paragraph{Bias elimination in the primal.}
For a fixed block, an admissible bias exists exactly when every score
at a positive-labelled query exceeds every score at a negative-labelled
query. If $a$ is positive and $b'$ negative, this is
\[
 (p_a-p_{b'})\sum_y
 \frac{u_{k,g}(\pi^{-1}y)c_y}
 {(1-p_a-p_y)(1-p_{b'}-p_y)}>0.                       \tag{3}
\]
The denominator factors are positive. Formula (3) follows by subtracting
the two simple reciprocals. Pure blocks impose no reduced inequalities.

\begin{lemma}[Positive kernel energy]
For distinct $x_1,\ldots,x_q$ in $[0,P]$, the matrix
$C_{ij}=1/(1-x_i-x_j)$ is positive definite. Consequently a nonzero
block vector $v$ has $v^TCv>0$, including when $\sum_i v_i=0$.
\end{lemma}
\begin{proof}
The quadratic form is
\[
 \int_0^\infty e^{-t}\left(\sum_i v_i e^{x_it}\right)^2\,dt.
\]
The integral converges since $2P<1$. If it is zero, the continuous
integrand is identically zero. Its first $q$ derivatives at zero give
a Vandermonde system, forcing every $v_i=0$.
\end{proof}

\begin{lemma}[Exact paired-leaf identity]
Fix a block $(k,g)$ and a lower output prefix $s<2^k$. There are unique
inputs $h_0,h_1$ in its two children whose final outputs have lower
prefix $s$. Their input bit $k$ is respectively zero and one. Put
$y_i=\pi(h_i)$, and let $L(h)$ denote the sum in (2) over levels below
$k$. For each strict ancestor block $a$ of $(k,g)$ its Haar sign
$u_a$ is the same at $h_0,h_1$. The difference of the two leaf equations
is
\[
 0=R_{k,g}(p_{y_0})+R_{k,g}(p_{y_1})
   +L(h_0)-L(h_1)
   +\sum_{a\text{ strict ancestor}}u_a
       \bigl(R_a(p_{y_0})-R_a(p_{y_1})\bigr).        \tag{4}
\]
\end{lemma}
\begin{proof}
Before layer $k$, routing is a permutation of the lower $k$ bits
within each fixed higher-bit string. The two inputs reaching the
layer-$k$ gate with prefix $s$ are therefore unique. Later layers
preserve bits below $k$. They have opposite signs in the block's Haar
function and equal signs in every strict ancestor. Subtracting the two
leaf equations proves (4).
\end{proof}

\begin{lemma}[Minimal-active-level estimate]
Suppose a dual certificate satisfies (2), and every block below level
$k\geq1$ is zero. Fix a block $g$ at level $k$, put $q=2^k$,
$M_g=\sum_s|v_{k,g,s}|$, and
$\Delta=\min_{s\ne t<q}|p_s-p_t|>0$.
Let $\delta_k=\sum_{j\geq k}\eta_j$, and set
\[
 \alpha_k=\frac{q^2\delta_k}{b^3\Delta^{2q-2}}.
\]
Then
\[
 M_g\leq\alpha_k\left(M_g+
                  \tfrac12\sum_{a\text{ strict ancestor}}M_a\right).
                                                               \tag{5}
\]
In particular, if $\alpha_k\leq1/2$, then
$M_g\leq\alpha_k\sum_aM_a$.
\end{lemma}
\begin{proof}
Every $p_{y_i}$ in (4) lies in $[p_s,p_s+\delta_k]$ and the two
points are at distance at most $\delta_k$.
On $[0,P]$ one has $|R_a'(x)|\leq M_a/b^2$.
Since $L=0$, (4) and the mean value theorem give
\[
 |R_{k,g}(p_s)|\leq\frac{\delta_k}{b^2}
       \left(M_g+\tfrac12\sum_aM_a\right).          \tag{6}
\]
The vector of these values equals $Cv_g$. The explicit inverse Cauchy
formula implies
\[
 |(C^{-1})_{ij}|
 =\frac{\prod_\ell(1-p_i-p_\ell)
             \prod_\ell(1-p_j-p_\ell)}
 {(1-p_i-p_j)\prod_{\ell\ne i}|p_i-p_\ell|
                  \prod_{\ell\ne j}|p_j-p_\ell|}
 \leq\frac1{b\Delta^{2q-2}}.
\]
Thus $\|C^{-1}x\|_1\leq q^2\|x\|_\infty/
(b\Delta^{2q-2})$. Applying this to (6) proves (5), and absorbing
the $\alpha_kM_g$ term proves the last assertion.
\end{proof}

\paragraph{The remaining compatibility gap.}
Estimate (5) is a one-level statement, not a triangular elimination of
all levels. After a lower block is bounded rather than exactly zero,
the term $L(h_0)-L(h_1)$ returns in (4). Later inverse-Cauchy norms
depend on the very small earlier digit gaps and can amplify that term.
Choosing each next digit small does not by itself establish that these
products are small. Likewise, the positive kernel energy does not yet
give a positive energy for the full masked leaf system.
An exact signed Schur-complement elimination, an appropriate weighted
energy identity, or a probabilistic control of these coupled equations
would be new work. Neither an all-target theorem nor a constant-density
theorem is asserted here.

\end{document}
